\textbf{Phần 1. Trắc nghiệm lựa chọn}
\setcounterpageref{ex}{0}

\begin{ex} %[5D7N42]
Phím nào trên máy tính cầm tay dùng để bật máy và xóa toàn bộ màn hình tính toán?
\choice
{\fbox{OFF}}
{\fbox{CE}}
{\fbox{=}}
{\True \fbox{ON/C}}
\loigiai{
Phím \fbox{ON/C} dùng để bật máy và xóa dữ liệu hiển thị trên màn hình.\\
\textbf{Đáp án của câu hỏi là D}
}
\end{ex}

\begin{ex} %[5D7N42]
Để thực hiện phép cộng trên máy tính cầm tay, ta sử dụng phím có kí hiệu nào?
\choice
{\fbox{--}}
{\fbox{$\times$}}
{\True \fbox{+}}
{\fbox{$\div$}}
\loigiai{
Phím phép cộng có kí hiệu \fbox{+}.\\
\textbf{Đáp án của câu hỏi là C}
}
\end{ex}

\begin{ex} %[5D7N42]
Để nhận kết quả của phép tính vừa nhập trên máy tính cầm tay, ta bấm phím nào?
\choice
{\fbox{ON}}
{\True \fbox{=}}
{\fbox{\%}}
{\fbox{CE}}
\loigiai{
Phím \fbox{=} dùng để tính và hiển thị kết quả.\\
\textbf{Đáp án của câu hỏi là B}
}
\end{ex}

\begin{ex} %[5D7H42]
Khi bấm lần lượt các phím: \fbox{2} \fbox{5} \fbox{+} \fbox{1} \fbox{8} \fbox{=}, màn hình máy tính sẽ hiển thị kết quả là
\choice
{\True $43$}
{$33$}
{$45$}
{$2518$}
\loigiai{
Kết quả: $25 + 18 = 43$.\\
\textbf{Đáp án của câu hỏi là A}
}
\end{ex}

\begin{ex} %[5D7H42]
Để tính $35\%$ của $400$ trên máy tính cầm tay có phím \%, ta có thể bấm lần lượt là
\choice
{\fbox{4} \fbox{0} \fbox{0} \fbox{+} \fbox{3} \fbox{5} \fbox{\%}}
{\fbox{4} \fbox{0} \fbox{0} \fbox{$\div$} \fbox{3} \fbox{5} \fbox{\%}}
{\True \fbox{4} \fbox{0} \fbox{0} \fbox{$\times$} \fbox{3} \fbox{5} \fbox{\%}}
{\fbox{3} \fbox{5} \fbox{+} \fbox{4} \fbox{0} \fbox{0} \fbox{=}}
\loigiai{
Cách tính $35\%$ của $400$: lấy $400 \times 35\%$.\\
\textbf{Đáp án của câu hỏi là C}
}
\end{ex}

\begin{ex} %[5D7H42]
Khi bấm các phím: \fbox{1} \fbox{.} \fbox{2} \fbox{5} \fbox{$\times$} \fbox{8} \fbox{=}, màn hình hiển thị kết quả là
\choice
{$100$}
{\True $10$}
{$1$}
{$0{,}1$}
\loigiai{
$1{,}25 \times 8 = 10$.\\
\textbf{Đáp án của câu hỏi là B}
}
\end{ex}

\begin{ex} %[5D7V42]
Để tính tỉ số phần trăm của $15$ và $60$ bằng máy tính cầm tay, ta bấm lần lượt các phím:
\choice
{\True \fbox{1} \fbox{5} \fbox{$\div$} \fbox{6} \fbox{0} \fbox{\%}}
{\fbox{1} \fbox{5} \fbox{$\times$} \fbox{6} \fbox{0} \fbox{=}}
{\fbox{6} \fbox{0} \fbox{$\div$} \fbox{1} \fbox{5} \fbox{=}}
{\fbox{1} \fbox{5} \fbox{+} \fbox{6} \fbox{0} \fbox{\%}}
\loigiai{
Để tìm tỉ số phần trăm: $15 : 60 \% = 25\%$.\\
\textbf{Đáp án của câu hỏi là A}
}
\end{ex}

\begin{ex} %[5D7V42]
Khi thực hiện phép tính $(45{,}6 + 18{,}4) : 4$ trên máy tính cầm tay, kết quả hiển thị trên màn hình là
\choice
{$14$}
{\True $16$}
{$18$}
{$20$}
\loigiai{
Ta có: $45{,}6 + 18{,}4 = 64$; $64 : 4 = 16$.\\
\textbf{Đáp án của câu hỏi là B}
}
\end{ex}

\begin{ex} %[5D7V42]
Một học sinh tính giá trị biểu thức $150 - 30 \times 4$ trên máy tính cầm tay không có chức năng thứ tự phép tính tự động (bấm lần lượt \fbox{1}\fbox{5}\fbox{0} \fbox{--} \fbox{3}\fbox{0} \fbox{$\times$} \fbox{4} \fbox{=}). Kết quả màn hình hiển thị là $480$. Kết quả đúng của biểu thức phải là
\choice
{$480$}
{$120$}
{$60$}
{\True $30$}
\loigiai{
Thứ tự phép tính: nhân trước trừ sau $\Rightarrow 150 - (30 \times 4) = 150 - 120 = 30$.\\
\textbf{Đáp án của câu hỏi là D}
}
\end{ex}

\begin{ex} %[5D7V42]
Khi bấm: \fbox{7} \fbox{2} \fbox{0} \fbox{$\times$} \fbox{1} \fbox{5} \fbox{\%} \fbox{=}, kết quả trên màn hình là
\choice
{$10{,}8$}
{$1\,080$}
{\True $108$}
{$54$}
\loigiai{
$720 \times 15\% = 108$.\\
\textbf{Đáp án của câu hỏi là C}
}
\end{ex}

\begin{ex} %[5D7V42]
Sử dụng máy tính cầm tay để tính kết quả biểu thức: $3{,}75 \times 4{,}8 - 2{,}5 \times 3{,}2$. Kết quả là
\choice
{$8$}
{$12$}
{$15$}
{\True $10$}
\loigiai{
$3{,}75 \times 4{,}8 = 18$; $2{,}5 \times 3{,}2 = 8$; $18 - 8 = 10$.\\
\textbf{Đáp án của câu hỏi là D}
}
\end{ex}

\begin{ex} %[5D7C42]
Sử dụng máy tính cầm tay tính giá trị của biểu thức: $\dfrac{12{,}6 \times 4{,}5 + 12{,}6 \times 5{,}5}{0{,}63 \times 4}$. Kết quả thu được là
\choice
{\True $50$}
{$25$}
{$100$}
{$200$}
\loigiai{
Tử số: $12{,}6 \times (4{,}5 + 5{,}5) = 12{,}6 \times 10 = 126$.\\
Mẫu số: $0{,}63 \times 4 = 2{,}52$.\\
Kết quả: $126 : 2{,}52 = 50$.\\
\textbf{Đáp án của câu hỏi là A}
}
\end{ex}

\textbf{Phần 2. Trắc nghiệm đúng sai}
\setcounterpageref{ex}{0}

\begin{ex} %[5D7H42]
Khi sử dụng máy tính cầm tay thông thường, xét tính đúng sai của các nhận định:
\choiceTF
{\True Phím \fbox{.} dùng để nhập dấu phẩy ngăn cách phần nguyên và phần thập phân.}
{Phím \fbox{OFF} dùng để xóa kết quả vừa tính mà không tắt máy.}
{\True Để tính $25 + 75$, ta bấm lần lượt: \fbox{2}\fbox{5} \fbox{+} \fbox{7}\fbox{5} \fbox{=}.}
{Bấm \fbox{5}\fbox{0} \fbox{$\div$} \fbox{0} \fbox{=} máy tính sẽ hiển thị số $0$.}
\loigiai{
\begin{itemize}
    \item Phím dấu chấm biểu thị dấu phân cách thập phân. Đúng.
    \item Phím OFF dùng để tắt nguồn máy tính. Sai.
    \item Thao tác cộng hai số nguyên: đúng.
    \item Phép chia cho số 0 là phép toán không hợp lệ, máy sẽ báo lỗi (Error). Sai.
\end{itemize}
\begin{itemize}
    \item a) \textbf{Đúng}.
    \item b) \textbf{Sai}.
    \item c) \textbf{Đúng}.
    \item d) \textbf{Sai}.
\end{itemize}
}
\end{ex}

\begin{ex} %[5D7V42]
Sử dụng máy tính cầm tay để tính các phép tính sau, xét tính đúng sai:
\choiceTF
{\True $125{,}6 \times 4{,}5 = 565{,}2$.}
{$\dfrac{3}{8}$ viết thành số thập phân bằng cách bấm \fbox{3} \fbox{$\div$} \fbox{8} \fbox{=} cho kết quả $0{,}38$.}
{\True $75\%$ của $640$ bằng cách bấm \fbox{6}\fbox{4}\fbox{0} \fbox{$\times$} \fbox{7}\fbox{5} \fbox{\%} cho kết quả $480$.}
{\True Phép chia $14 : 25$ hiển thị kết quả là $0{,}56$.}
\loigiai{
\begin{itemize}
    \item $125{,}6 \times 4{,}5 = 565{,}2$. Đúng.
    \item $3 : 8 = 0{,}375 \neq 0{,}38$. Sai.
    \item $640 \times 75\% = 480$. Đúng.
    \item $14 : 25 = 0{,}56$. Đúng.
\end{itemize}
\begin{itemize}
    \item a) \textbf{Đúng}.
    \item b) \textbf{Sai}.
    \item c) \textbf{Đúng}.
    \item d) \textbf{Đúng}.
\end{itemize}
}
\end{ex}

\textbf{Phần 3. Trả lời ngắn}
\setcounterpageref{ex}{0}

\begin{ex} %[5D7H42]
Dùng máy tính cầm tay thực hiện phép tính: $456{,}8 + 289{,}75$. (Ghi kết quả dưới dạng số thập phân dùng dấu phẩy).
\shortans{746,55}
\loigiai{
\textbf{Đáp số:} $746{,}55$\\
$456{,}8 + 289{,}75 = 746{,}55$.
}
\end{ex}

\begin{ex} %[5D7H42]
Dùng máy tính cầm tay tính kết quả: $85{,}4 \times 6{,}5$.
\shortans{555,1}
\loigiai{
\textbf{Đáp số:} $555{,}1$\\
$85{,}4 \times 6{,}5 = 555{,}1$.
}
\end{ex}

\begin{ex} %[5D7V42]
Dùng máy tính cầm tay tính kết quả: $148{,}8 : 4{,}8$.
\shortans{31}
\loigiai{
\textbf{Đáp số:} $31$\\
$148{,}8 : 4{,}8 = 31$.
}
\end{ex}

\begin{ex} %[5D7V42]
Sử dụng máy tính cầm tay để tìm $18\%$ của số $1\,250$.
\shortans{225}
\loigiai{
\textbf{Đáp số:} $225$\\
$1\,250 \times 18\% = 225$.
}
\end{ex}

\begin{ex} %[5D7V42]
Dùng máy tính cầm tay tính giá trị biểu thức: $(48{,}6 + 31{,}4) \times 2{,}5 - 50$.
\shortans{150}
\loigiai{
\textbf{Đáp số:} $150$\\
$80 \times 2{,}5 - 50 = 200 - 50 = 150$.
}
\end{ex}

\begin{ex} %[5D7C42]
Tính giá trị biểu thức bằng máy tính cầm tay: $15{,}6 \times 3{,}4 + 15{,}6 \times 6{,}6 + 44$.
\shortans{200}
\loigiai{
\textbf{Đáp số:} $200$\\
$15{,}6 \times (3{,}4 + 6{,}6) + 44 = 156 + 44 = 200$.
}
\end{ex}
